\chapter{Tableaux, tri et enregistrements en C}

\begin{synthese}[Au programme]
Ranger plusieurs valeurs sous un même nom grâce aux \textbf{tableaux}, les
parcourir avec une boucle, calculer somme, moyenne, minimum et maximum, rechercher
une valeur, puis \textbf{trier par sélection}. On découvre aussi les
\textbf{matrices} (tableaux à deux dimensions) et les \textbf{enregistrements}
(\code{struct}) pour décrire une fiche élève. C'est le cœur de la manipulation de
données en langage C.
\end{synthese}

\section{Les tableaux à une dimension}

Jusqu'ici, chaque valeur tenait dans une variable. Mais comment stocker les
\textbf{60 notes} d'une classe de Terminale~TI ? Déclarer \code{n1}, \code{n2},
\dots, \code{n60} serait absurde. Le \textbf{tableau} résout ce problème : un seul
nom regroupe plusieurs cases de \textbf{même type}, numérotées par un \textbf{indice}.

\begin{definition}
Un \textbf{tableau} (\emph{array}) est une suite de cases mémoire contiguës, de
même type, désignées par un nom unique. On accède à chaque case par son
\textbf{indice} (sa position), qui commence à \textbf{0} en langage C.
\end{definition}

\subsection{Déclaration}

\begin{syntaxe}
\code{type nom[taille];}\quad — par exemple \code{int t[10];}
\end{syntaxe}

La déclaration \code{int t[10];} réserve \textbf{10 cases} d'entiers, d'indices
\code{0} à \code{9}. La \textbf{taille} doit être une constante connue à la
compilation.

\begin{codec}
int t[10];          // 10 entiers : t[0], t[1], ..., t[9]
float notes[60];    // 60 notes décimales
char nom[30];       // 30 caractères (une chaîne)
\end{codec}

\begin{attention}
Un tableau de taille \code{10} possède les indices \textbf{0 à 9}. La case
\code{t[10]} \textbf{n'existe pas} : y accéder est une erreur classique
(\emph{débordement}) qui plante le programme ou corrompt la mémoire.
\end{attention}

\subsection{Initialisation et accès}

On peut donner les valeurs dès la déclaration, entre accolades~:

\begin{codec}
int t[5] = {12, 7, 25, 3, 18};   // les 5 cases sont remplies
int u[5] = {0};                  // toutes les cases à 0
int v[] = {4, 8, 15};            // taille déduite : 3 cases
\end{codec}

On lit ou modifie une case avec \code{nom[indice]}~:

\begin{codec}
t[0] = 20;              // on change la 1re case
int x = t[2];           // x reçoit la valeur de la 3e case (25)
t[4] = t[4] + 1;        // on incrémente la dernière case
\end{codec}

\begin{remarque}
L'indice peut être une variable~: \code{t[i]}. C'est ce qui permet de parcourir
tout le tableau avec une boucle \code{for} en faisant varier \code{i}.
\end{remarque}

\subsection{Saisie et affichage par une boucle}

Pour remplir ou afficher un tableau, on utilise une boucle \code{for} dont le
compteur sert d'indice.

\begin{codec}
// Programme en langage C
#include <stdio.h>

int main() {
    int notes[5];
    int i;

    // SAISIE des 5 notes
    for (i = 0; i < 5; i++) {
        printf("Note numero %d : ", i + 1);
        scanf("%d", &notes[i]);     // &notes[i] : adresse de la case
    }

    // AFFICHAGE des 5 notes
    printf("Vous avez saisi : ");
    for (i = 0; i < 5; i++) {
        printf("%d ", notes[i]);
    }
    printf("\n");

    return 0;
}
\end{codec}

\fig{La boucle parcourt les indices 0 à 4 : un seul code pour toutes les cases.}

\begin{attention}
Dans \code{scanf}, on met le \code{\&} devant la case~: \code{scanf("\%d", \&notes[i])}.
Il indique l'\textbf{adresse} où ranger la valeur lue.
\end{attention}

\subsection{Passer un tableau à une fonction}

Quand on transmet un tableau à une fonction, le C passe en réalité l'\textbf{adresse}
de sa première case. La fonction travaille donc \textbf{directement} sur le tableau
original~: toute modification est conservée. On doit aussi passer la \textbf{taille},
car la fonction ne la connaît pas.

\begin{codec}
// Programme en langage C
#include <stdio.h>

// Affiche les n premieres cases du tableau t
void afficher(int t[], int n) {
    int i;
    for (i = 0; i < n; i++) {
        printf("%d ", t[i]);
    }
    printf("\n");
}

int main() {
    int tab[4] = {10, 20, 30, 40};
    afficher(tab, 4);     // affiche : 10 20 30 40
    return 0;
}
\end{codec}

\begin{remarque}
Dans l'en-tête \code{void afficher(int t[], int n)}, les crochets \code{[]} restent
\textbf{vides}~: la fonction accepte un tableau de n'importe quelle longueur, et
\code{n} précise combien de cases utiliser.
\end{remarque}

\section{Algorithmes classiques sur un tableau}

Une poignée de schémas reviennent sans cesse. Apprends-les par cœur~: tous reposent
sur une boucle qui parcourt le tableau.

\subsection{Somme et moyenne}

On utilise un \textbf{accumulateur} initialisé à \code{0}, auquel on ajoute chaque
case.

\begin{codec}
// Programme en langage C
#include <stdio.h>

int main() {
    float notes[5] = {12.0, 15.5, 8.0, 14.0, 11.5};
    int i;
    float somme = 0.0;

    for (i = 0; i < 5; i++) {
        somme = somme + notes[i];   // on accumule
    }
    float moyenne = somme / 5;      // division par le nombre de notes

    printf("Somme   : %.2f\n", somme);
    printf("Moyenne : %.2f\n", moyenne);
    return 0;
}
\end{codec}

\fig{Sortie : Somme : 61.00 \quad Moyenne : 12.20}

\subsection{Minimum et maximum}

On suppose d'abord que la \textbf{première case} est le maximum, puis on la compare
à toutes les autres et on garde la plus grande.

\begin{methode}
\textbf{Chercher le maximum.} (1)~\code{max} reçoit \code{t[0]}~; (2)~pour chaque
case suivante, si \code{t[i] > max} alors \code{max} reçoit \code{t[i]}~; (3)~à la
fin, \code{max} contient la plus grande valeur. (Pour le minimum, on remplace
\code{>} par \code{<}.)
\end{methode}

\begin{codec}
// Programme en langage C
#include <stdio.h>

int main() {
    int t[6] = {14, 9, 23, 7, 31, 18};
    int i;
    int max = t[0];     // hypothese : la 1re case est le max

    for (i = 1; i < 6; i++) {       // on commence a l'indice 1
        if (t[i] > max) {
            max = t[i];
        }
    }
    printf("Le maximum est %d\n", max);   // affiche 31
    return 0;
}
\end{codec}

\begin{attention}
On initialise \code{max} avec \code{t[0]}, \textbf{jamais} avec \code{0}~: si toutes
les valeurs étaient négatives, partir de \code{0} donnerait un résultat faux.
\end{attention}

\subsection{Recherche séquentielle d'une valeur}

Pour savoir si une valeur \code{cherchee} se trouve dans le tableau, on parcourt
les cases une à une~: c'est la \textbf{recherche séquentielle} (ou linéaire).

\begin{codec}
// Programme en langage C
#include <stdio.h>

int main() {
    int t[6] = {14, 9, 23, 7, 31, 18};
    int cherchee = 23;
    int i;
    int trouve = 0;     // 0 = pas encore trouve (faux)
    int position = -1;

    for (i = 0; i < 6; i++) {
        if (t[i] == cherchee) {
            trouve = 1;         // on l'a trouvee (vrai)
            position = i;       // on retient sa position
        }
    }

    if (trouve == 1) {
        printf("Valeur trouvee a l'indice %d\n", position);
    } else {
        printf("Valeur absente du tableau\n");
    }
    return 0;
}
\end{codec}

\fig{Le drapeau \code{trouve} (0/1) mémorise le résultat de la recherche.}

\section{Le tri par sélection}

\textbf{Trier} un tableau, c'est ranger ses valeurs dans l'ordre croissant (ou
décroissant). Le \textbf{tri par sélection} est la méthode la plus simple à
comprendre et figure au programme de Terminale~TI.

\subsection{Le principe}

\begin{definition}
Le \textbf{tri par sélection} consiste à~: chercher le plus petit élément du
tableau et le placer en première position~; puis chercher le plus petit parmi
\textbf{ce qui reste} et le placer en deuxième position~; et ainsi de suite jusqu'à
la fin.
\end{definition}

À chaque tour, on \textbf{sélectionne} le minimum de la partie non triée et on
l'\textbf{échange} avec la première case de cette partie.

\begin{methode}
\textbf{Tri par sélection (pseudo-algorithme).}
\begin{enumerate}[label=\arabic*)]
  \item Pour \code{i} allant de \code{0} à \code{n-2}~:
  \item \quad poser \code{iMin = i} (position du plus petit présumé)~;
  \item \quad pour \code{j} allant de \code{i+1} à \code{n-1}~: si \code{t[j] < t[iMin]} alors \code{iMin = j}~;
  \item \quad échanger \code{t[i]} et \code{t[iMin]}.
\end{enumerate}
\end{methode}

\subsection{Trace sur un exemple}

Trions le tableau \code{[5, 3, 8, 1]} ($n = 4$). Chaque ligne montre l'état après
l'échange du tour~$i$ ; la partie \textbf{triée} est à gauche du trait.

\begin{center}\small
\begin{tabular}{@{}clll@{}}
\toprule
\textbf{Tour} & \textbf{Plus petit trouvé} & \textbf{Échange} & \textbf{Tableau}\\
\midrule
Départ & --- & --- & 5\ \ 3\ \ 8\ \ 1\\
$i=0$ & 1 (indice 3) & t[0] $\leftrightarrow$ t[3] & \textbf{1} $|$ 3\ \ 8\ \ 5\\
$i=1$ & 3 (indice 1) & aucun (déjà en place) & 1\ \ \textbf{3} $|$ 8\ \ 5\\
$i=2$ & 5 (indice 3) & t[2] $\leftrightarrow$ t[3] & 1\ \ 3\ \ \textbf{5} $|$ 8\\
Fin & --- & --- & 1\ \ 3\ \ 5\ \ 8\\
\bottomrule
\end{tabular}
\end{center}

\fig{En 3 tours (n-1), le tableau est trié dans l'ordre croissant.}

\subsection{Code C complet}

\begin{codec}
// Programme en langage C
#include <stdio.h>

int main() {
    int t[6] = {5, 3, 8, 1, 9, 2};
    int n = 6;
    int i, j, iMin, temp;

    // TRI PAR SELECTION
    for (i = 0; i < n - 1; i++) {
        iMin = i;                       // plus petit presume : la case i
        for (j = i + 1; j < n; j++) {
            if (t[j] < t[iMin]) {
                iMin = j;               // on a trouve plus petit
            }
        }
        // echange de t[i] et t[iMin] via une variable temporaire
        temp = t[i];
        t[i] = t[iMin];
        t[iMin] = temp;
    }

    // affichage du tableau trie
    printf("Tableau trie : ");
    for (i = 0; i < n; i++) {
        printf("%d ", t[i]);
    }
    printf("\n");
    return 0;
}
\end{codec}

\fig{Sortie : Tableau trie : 1 2 3 5 8 9}

\begin{attention}
L'\textbf{échange} de deux cases passe \textbf{obligatoirement} par une variable
temporaire (\code{temp}). Écrire directement \code{t[i] = t[iMin];} puis
\code{t[iMin] = t[i];} écrase la première valeur et perd une donnée.
\end{attention}

\section{Les tableaux à deux dimensions (matrices)}

Un tableau à \textbf{deux dimensions} se représente comme un \textbf{tableau de
lignes et de colonnes}, exactement comme un tableau de notes par élève et par
matière.

\begin{syntaxe}
\code{type nom[lignes][colonnes];}\quad — par exemple \code{int m[3][4];}
\end{syntaxe}

\code{int m[3][4];} déclare une matrice de \textbf{3 lignes} et \textbf{4 colonnes},
soit 12 cases. On accède à une case par \code{m[i][j]} (ligne \code{i}, colonne
\code{j}).

\begin{codec}
// Programme en langage C
#include <stdio.h>

int main() {
    int m[2][3] = { {1, 2, 3},
                    {4, 5, 6} };   // 2 lignes, 3 colonnes
    int i, j;

    // PARCOURS DOUBLE : une boucle par dimension
    for (i = 0; i < 2; i++) {           // pour chaque ligne
        for (j = 0; j < 3; j++) {       // pour chaque colonne
            printf("%d\t", m[i][j]);
        }
        printf("\n");       // retour a la ligne apres chaque ligne
    }
    return 0;
}
\end{codec}

\fig{Le parcours d'une matrice utilise deux boucles imbriquées (lignes, colonnes).}

\begin{remarque}
La boucle \textbf{extérieure} parcourt les lignes (\code{i}), la boucle
\textbf{intérieure} parcourt les colonnes (\code{j}). On affiche un retour à la
ligne (\code{\\n}) après chaque ligne pour obtenir un beau tableau.
\end{remarque}

\section{Les enregistrements (struct)}

Un tableau regroupe des valeurs de \textbf{même type}. Mais une fiche élève réunit
un \textbf{nom} (texte), un \textbf{âge} (entier) et une \textbf{moyenne} (décimal)~:
des types différents. L'\textbf{enregistrement} (\code{struct}) permet de les
rassembler sous un seul nom.

\begin{definition}
Un \textbf{enregistrement} (mot-clé \code{struct}) est un type composé qui regroupe
plusieurs variables, appelées \textbf{champs}, pouvant être de types différents.
\end{definition}

\subsection{Définir une structure}

\begin{codec}
struct Eleve {
    char nom[30];       // champ texte
    int age;            // champ entier
    float moyenne;      // champ decimal
};                      // ne pas oublier le point-virgule final
\end{codec}

\begin{attention}
La définition d'une \code{struct} se termine par un \textbf{point-virgule} après
l'accolade fermante. C'est une erreur de débutant très fréquente de l'oublier.
\end{attention}

\subsection{Déclarer une variable et accéder aux champs}

On accède à un champ avec le \textbf{point} (\code{.})~: \code{variable.champ}.

\begin{codec}
// Programme en langage C
#include <stdio.h>
#include <string.h>

struct Eleve {
    char nom[30];
    int age;
    float moyenne;
};

int main() {
    struct Eleve e1;            // une variable de type struct Eleve

    strcpy(e1.nom, "Awa");      // copie un texte dans un champ char[]
    e1.age = 17;                // affectation directe pour les autres
    e1.moyenne = 14.5;

    printf("Nom     : %s\n", e1.nom);
    printf("Age     : %d ans\n", e1.age);
    printf("Moyenne : %.2f\n", e1.moyenne);
    return 0;
}
\end{codec}

\begin{remarque}
Un champ \code{char nom[30]} est un tableau de caractères~: on ne peut pas écrire
\code{e1.nom = "Awa";}. On copie le texte avec \code{strcpy} (de la bibliothèque
\code{<string.h>}).
\end{remarque}

\subsection{Un tableau de structures}

On peut créer un \textbf{tableau de \code{struct}} pour gérer toute une classe~:
chaque case est une fiche complète.

\begin{codec}
// Programme en langage C
#include <stdio.h>
#include <string.h>

struct Eleve {
    char nom[30];
    int age;
    float moyenne;
};

int main() {
    struct Eleve classe[2];     // tableau de 2 eleves

    strcpy(classe[0].nom, "Awa");
    classe[0].age = 17;
    classe[0].moyenne = 14.5;

    strcpy(classe[1].nom, "Bello");
    classe[1].age = 18;
    classe[1].moyenne = 11.0;

    int i;
    for (i = 0; i < 2; i++) {
        printf("%s (%d ans) : %.2f\n",
               classe[i].nom, classe[i].age, classe[i].moyenne);
    }
    return 0;
}
\end{codec}

\fig{Chaque case \code{classe[i]} est un enregistrement complet (nom, âge, moyenne).}

\section{Exemple complet : la fiche élève}

Réunissons tout~: saisir une classe, l'afficher, puis trouver le \textbf{meilleur
élève} (plus forte moyenne).

\begin{codec}
// Programme en langage C
#include <stdio.h>

struct Eleve {
    char nom[30];
    int age;
    float moyenne;
};

int main() {
    struct Eleve classe[3];
    int n = 3;
    int i;

    // SAISIE
    for (i = 0; i < n; i++) {
        printf("Eleve %d - nom : ", i + 1);
        scanf("%s", classe[i].nom);     // pas de & : nom est deja un tableau
        printf("Eleve %d - age : ", i + 1);
        scanf("%d", &classe[i].age);
        printf("Eleve %d - moyenne : ", i + 1);
        scanf("%f", &classe[i].moyenne);
    }

    // AFFICHAGE
    printf("\n--- Liste de la classe ---\n");
    for (i = 0; i < n; i++) {
        printf("%s : %.2f\n", classe[i].nom, classe[i].moyenne);
    }

    // MEILLEUR ELEVE (plus forte moyenne)
    int iMax = 0;
    for (i = 1; i < n; i++) {
        if (classe[i].moyenne > classe[iMax].moyenne) {
            iMax = i;
        }
    }
    printf("\nMajor de la classe : %s (%.2f)\n",
           classe[iMax].nom, classe[iMax].moyenne);
    return 0;
}
\end{codec}

\fig{On retrouve le schéma du \emph{maximum}, appliqué au champ \code{moyenne}.}

\begin{remarque}
Pour lire un nom avec \code{scanf("\%s", classe[i].nom)}, on ne met \textbf{pas} de
\code{\&}~: \code{classe[i].nom} est un tableau, c'est déjà une adresse. En
revanche, \code{age} et \code{moyenne} en réclament un.
\end{remarque}

\section{Exercices}

\rubrique{Application}

\exo{Somme et moyenne}\diff{1}
Un tableau \code{int notes[5] = \{12, 8, 15, 10, 14\};} contient 5 notes. Écris un
programme qui calcule et affiche leur somme et leur moyenne.

\exo{Minimum et maximum}\diff{1}
À partir du tableau \code{int t[6] = \{17, 4, 23, 9, 1, 30\};}, affiche la plus
petite et la plus grande valeur.

\exo{Compter les admis}\diff{2}
Un tableau \code{float moy[8]} contient les moyennes de 8 élèves d'un lycée de
Bafoussam. Écris un programme qui compte et affiche le nombre d'élèves
\textbf{admis} (moyenne supérieure ou égale à 10).

\exo{Recherche d'une valeur}\diff{2}
On donne un tableau \code{int t[7] = \{5, 12, 8, 20, 3, 17, 9\};} et un entier
\code{cherchee = 20}. Indique si cette valeur est présente, et à quelle position.

\rubrique{Entraînement}

\exo{Inverser un tableau}\diff{2}
Écris un programme qui inverse l'ordre des éléments d'un tableau \code{t} de 6
entiers (par exemple les valeurs 1 à 6)~: la première case échange avec la
dernière, la deuxième avec l'avant-dernière, etc. Affiche ensuite le résultat.

\exo{Tri par sélection}\diff{3}
Saisis 5 entiers dans un tableau, trie-le par \textbf{sélection} dans l'ordre
croissant, puis affiche le tableau trié.

\exo{Gestion de produits}\diff{3}
Une boutique de Douala gère ses articles avec une structure \code{Produit}
possédant les champs \code{nom} (texte), \code{quantite} (entier) et \code{prix}
(décimal). Saisis 3~produits, affiche-les, puis indique le produit le
\textbf{plus cher} (en francs CFA).

\exo{Afficher une matrice}\diff{2}
Déclare et initialise une matrice \code{int m[3][3]} avec des valeurs de ton choix,
puis affiche-la sous forme de tableau (lignes et colonnes alignées).

\section{Corrigés}

\corrige{de l'exercice 1}
\begin{codec}
// Programme en langage C
#include <stdio.h>

int main() {
    int notes[5] = {12, 8, 15, 10, 14};
    int i, somme = 0;

    for (i = 0; i < 5; i++) {
        somme = somme + notes[i];
    }
    float moyenne = (float) somme / 5;    // (float) pour un resultat decimal

    printf("Somme   : %d\n", somme);
    printf("Moyenne : %.2f\n", moyenne);
    return 0;
}
\end{codec}
La somme vaut \code{59} et la moyenne \code{11.80}. Le \code{(float)} force la
division décimale (sinon \code{59 / 5} donnerait \code{11} en entier).

\corrige{de l'exercice 2}
\begin{codec}
// Programme en langage C
#include <stdio.h>

int main() {
    int t[6] = {17, 4, 23, 9, 1, 30};
    int i;
    int min = t[0], max = t[0];

    for (i = 1; i < 6; i++) {
        if (t[i] < min) min = t[i];
        if (t[i] > max) max = t[i];
    }
    printf("Minimum : %d\n", min);   // 1
    printf("Maximum : %d\n", max);   // 30
    return 0;
}
\end{codec}
Une seule boucle suffit pour trouver le minimum \textbf{et} le maximum, en
initialisant les deux avec \code{t[0]}.

\corrige{de l'exercice 3}
\begin{codec}
// Programme en langage C
#include <stdio.h>

int main() {
    float moy[8] = {12.0, 8.5, 15.0, 9.0, 10.0, 7.5, 13.5, 11.0};
    int i, admis = 0;

    for (i = 0; i < 8; i++) {
        if (moy[i] >= 10.0) {
            admis = admis + 1;      // on compte
        }
    }
    printf("Nombre d'admis : %d sur 8\n", admis);   // 5
    return 0;
}
\end{codec}
On utilise un \textbf{compteur} (\code{admis}) incrémenté chaque fois que la
condition \code{moy[i] >= 10.0} est vraie.

\corrige{de l'exercice 4}
\begin{codec}
// Programme en langage C
#include <stdio.h>

int main() {
    int t[7] = {5, 12, 8, 20, 3, 17, 9};
    int cherchee = 20;
    int i, position = -1;

    for (i = 0; i < 7; i++) {
        if (t[i] == cherchee) {
            position = i;       // on memorise l'indice
        }
    }

    if (position != -1) {
        printf("Valeur %d trouvee a l'indice %d\n", cherchee, position);
    } else {
        printf("Valeur %d absente\n", cherchee);
    }
    return 0;
}
\end{codec}
Ici \code{20} est à l'indice \code{3}. La variable \code{position} vaut \code{-1}
tant que rien n'est trouvé, ce qui sert de signal d'absence.

\corrige{de l'exercice 5}
\begin{codec}
// Programme en langage C
#include <stdio.h>

int main() {
    int t[6] = {1, 2, 3, 4, 5, 6};
    int n = 6, i, temp;

    // on echange t[i] avec t[n-1-i], jusqu'au milieu
    for (i = 0; i < n / 2; i++) {
        temp = t[i];
        t[i] = t[n - 1 - i];
        t[n - 1 - i] = temp;
    }

    for (i = 0; i < n; i++) {
        printf("%d ", t[i]);    // affiche : 6 5 4 3 2 1
    }
    printf("\n");
    return 0;
}
\end{codec}
On ne parcourt que la \textbf{moitié} du tableau (\code{i < n/2})~: chaque tour
échange une case du début avec la case symétrique de la fin.

\corrige{de l'exercice 6}
\begin{codec}
// Programme en langage C
#include <stdio.h>

int main() {
    int t[5];
    int i, j, iMin, temp;

    for (i = 0; i < 5; i++) {
        printf("Entier %d : ", i + 1);
        scanf("%d", &t[i]);
    }

    // tri par selection
    for (i = 0; i < 4; i++) {           // jusqu'a n-2 (soit 3)
        iMin = i;
        for (j = i + 1; j < 5; j++) {
            if (t[j] < t[iMin]) iMin = j;
        }
        temp = t[i];
        t[i] = t[iMin];
        t[iMin] = temp;
    }

    printf("Tableau trie : ");
    for (i = 0; i < 5; i++) printf("%d ", t[i]);
    printf("\n");
    return 0;
}
\end{codec}

\corrige{de l'exercice 7}
\begin{codec}
// Programme en langage C
#include <stdio.h>

struct Produit {
    char nom[30];
    int quantite;
    float prix;
};

int main() {
    struct Produit stock[3];
    int i, iCher = 0;

    for (i = 0; i < 3; i++) {
        printf("Produit %d - nom : ", i + 1);
        scanf("%s", stock[i].nom);
        printf("Produit %d - quantite : ", i + 1);
        scanf("%d", &stock[i].quantite);
        printf("Produit %d - prix (FCFA) : ", i + 1);
        scanf("%f", &stock[i].prix);
    }

    printf("\n--- Inventaire ---\n");
    for (i = 0; i < 3; i++) {
        printf("%s : %d unites a %.0f FCFA\n",
               stock[i].nom, stock[i].quantite, stock[i].prix);
    }

    // produit le plus cher
    for (i = 1; i < 3; i++) {
        if (stock[i].prix > stock[iCher].prix) iCher = i;
    }
    printf("\nProduit le plus cher : %s (%.0f FCFA)\n",
           stock[iCher].nom, stock[iCher].prix);
    return 0;
}
\end{codec}

\corrige{de l'exercice 8}
\begin{codec}
// Programme en langage C
#include <stdio.h>

int main() {
    int m[3][3] = { {1, 2, 3},
                    {4, 5, 6},
                    {7, 8, 9} };
    int i, j;

    for (i = 0; i < 3; i++) {
        for (j = 0; j < 3; j++) {
            printf("%d\t", m[i][j]);
        }
        printf("\n");
    }
    return 0;
}
\end{codec}
Les deux boucles imbriquées parcourent lignes (\code{i}) puis colonnes (\code{j})~;
le \code{\\n} après la boucle interne sépare les lignes.

\begin{aretenir}
Un \textbf{tableau} regroupe des valeurs de même type, indexées à partir de
\textbf{0}, qu'on parcourt avec une boucle \code{for}. Les schémas
\textbf{somme/moyenne}, \textbf{min/max} et \textbf{recherche} sont à connaître par
cœur. Le \textbf{tri par sélection} place à chaque tour le plus petit élément
restant, via un \textbf{échange} avec variable temporaire. Une \code{struct}
rassemble des champs de types différents (fiche élève)~; on y accède avec le point
(\code{.}) et on copie un nom avec \code{strcpy}.
\end{aretenir}
